\documentclass{beamer}

% --- Theme and layout ---
\usetheme{Darmstadt}
\usecolortheme{default}
\usefonttheme{default}
\usefonttheme[onlymath]{serif}
\setbeamertemplate{navigation symbols}{}
\setbeamertemplate{footline}{}
\setbeamertemplate{caption}[numbered]
\setbeamertemplate{blocks}[rounded][shadow=true]
\usepackage{tikz}
\usetikzlibrary{arrows.meta,positioning}

\definecolor{raceblue}{RGB}{35,110,180}
\definecolor{othergray}{RGB}{215,220,225}

% Slightly reduce spacing around displayed equations
\AtBeginDocument{
  \setlength{\abovedisplayskip}{5pt}
  \setlength{\belowdisplayskip}{5pt}
  \setlength{\abovedisplayshortskip}{3pt}
  \setlength{\belowdisplayshortskip}{3pt}
}



\usepackage[T1]{fontenc}
\usepackage{lmodern}
\usepackage{amsmath,amssymb}
\usepackage{tikz}
\usepackage{pgfplots}
\usetikzlibrary{arrows.meta,positioning}
\pgfplotsset{compat=1.18}

\definecolor{raceblue}{RGB}{35,110,180}
\definecolor{othergray}{RGB}{215,220,225}

\setbeamercolor{structure}{fg=raceblue}
\setbeamertemplate{navigation symbols}{}

% Left and right margins
\setbeamersize{
    text margin left=0.8cm,
    text margin right=0.8cm
}

% Space above and below each slide title
\addtobeamertemplate{frametitle}
    {\vspace*{0.20cm}}
    {\vspace*{0.10cm}}

% Reserve bottom space and keep page numbers away from the edge
\setbeamertemplate{footline}{%
    \leavevmode
    \hbox{%
        \begin{beamercolorbox}[
            wd=\paperwidth,
            ht=0.45cm,
            dp=0.25cm,
            right
        ]{normal text}%
            \usebeamerfont{page number in head/foot}%
            \insertframenumber\hspace*{0.8cm}%
        \end{beamercolorbox}%
    }%
}

% Two equally populated neighborhoods.
% Argument 1: number of blue dots in the upper neighborhood.
% Argument 2: number of blue dots in the lower neighborhood.
% Arguments 3 and 4: neighborhood group shares.
\newcommand{\citydiagram}[4]{%
\begin{tikzpicture}[x=0.43cm,y=0.43cm]
    \draw[rounded corners,gray]
        (-0.6,-0.6) rectangle (4.6,1.6);
    \draw[rounded corners,gray]
        (-0.6,-3.6) rectangle (4.6,-1.4);

    \foreach \x in {0,...,4}{
        \foreach \y in {0,1}{
            \fill[othergray] (\x,\y) circle (0.19);
            \fill[othergray] (\x,\y-3) circle (0.19);
        }
    }

    \ifnum#1>0
        \foreach \i in {1,...,#1}{
            \fill[raceblue] ({\i-1},1) circle (0.19);
        }
    \fi
    \ifnum#2>0
        \foreach \i in {1,...,#2}{
            \fill[raceblue] ({\i-1},-2) circle (0.19);
        }
    \fi

    \node[right,font=\scriptsize] at (4.8,0.5)
        {$r_1=#3$};
    \node[right,font=\scriptsize] at (4.8,-2.5)
        {$r_2=#4$};

    \draw[<->,>=Stealth,gray,thick]
        (2,-0.65) -- (2,-1.35);
\end{tikzpicture}%
}
% --- Color palette ---
\definecolor{MyNavy}{RGB}{25,58,86}
\definecolor{MyTeal}{RGB}{43,119,125}
\definecolor{MyLightBlue}{RGB}{232,239,244}

\setbeamercolor{normal text}{fg=black,bg=white}
\setbeamercolor{structure}{fg=MyNavy}
\setbeamercolor{frametitle}{fg=white,bg=MyNavy}
\setbeamercolor{title}{fg=white,bg=MyNavy}
\setbeamercolor{titlelike}{fg=white,bg=MyNavy}
\setbeamercolor{section in head/foot}{fg=white,bg=MyNavy}
\setbeamercolor{subsection in head/foot}{fg=white,bg=MyTeal}
\setbeamercolor{block title}{fg=white,bg=MyNavy}
\setbeamercolor{block body}{fg=black,bg=MyLightBlue}
\setbeamercolor{alerted text}{fg=MyTeal}
\setbeamercolor{author}{fg=black}
\setbeamercolor{date}{fg=black}
\setbeamercolor{institute}{fg=black}

% --- Packages ---
\usepackage{amsmath,amssymb,amsbsy}
\usepackage{graphicx}
\usepackage{url}
\usepackage{biblatex}
\usepackage{algorithm}
\usepackage[noend]{algpseudocode}
\usepackage{siunitx}
\usepackage{tikz}
\usepackage{appendixnumberbeamer}


\begin{document}
%% Replace the entire supplied frame excerpt with these frames.
%% Keep your existing preamble unchanged.

\begin{frame}{Proposition 5.1: Why Monotonicity Holds (I)}
\small

Suppose $0\leq B\leq C$ entrywise.

Because both matrices are nonnegative,
multiplication preserves the inequalities:
\[
0\leq B^k\leq C^k
\qquad(k\geq1).
\]

Using the infinity norm,
\[
\|B^k\|_\infty\leq\|C^k\|_\infty.
\]
\end{frame}

\begin{frame}{Proposition 5.1: Why Monotonicity Holds (II)}
\small

We have $\|B^k\|_\infty\leq\|C^k\|_\infty$ for every $k\geq1$.

By Gelfand's formula,
\[
\rho(B)=\lim_{k\to\infty}
\|B^k\|_\infty^{1/k}.
\]

Consequently,
\[
\boxed{\rho(B)\leq\rho(C).}
\]

\begin{block}{Application}
Taking $B=cJ$ and $C=A$ proves
that $\rho(A)>0$.
\end{block}
\end{frame}

\begin{frame}{Proposition 5.1: Strict Dominance (I)}
\small

Normalize $A$ by its spectral radius:
\[
M=\frac{A}{\rho(A)},\qquad \rho(M)=1.
\]

Choose an eigenvalue $\mu$ of $M$ satisfying
$|\mu|=1$.

We want to show that $\mu=1$.

Suppose instead that $\mu\neq1$.
Choose an integer $m\geq1$ such that
\[
\operatorname{Re}(\mu^m)<0.
\]
\end{frame}

\begin{frame}{Proposition 5.1: Strict Dominance (II)}
\small

Recall that $|\mu|=1$ and $\operatorname{Re}(\mu^m)<0$.

Since $M>0$, we also have $M^m>0$.

For sufficiently small $\varepsilon>0$,
define
\[
T=M^m-\varepsilon I\geq0.
\]

\begin{block}{Key Observation}
Because $\mu^m$ is an eigenvalue of $M^m$,
the number $\mu^m-\varepsilon$ is an
eigenvalue of $T$.
\end{block}
\end{frame}

\begin{frame}{Proposition 5.1: Strict Dominance (III)}
\small

From the previous slide,
$\mu^m-\varepsilon$ is an eigenvalue of $T$.

Since $|\mu|=1$ and
$\operatorname{Re}(\mu^m)<0$,
\[
\begin{aligned}
|\mu^m-\varepsilon|^2
&=|\mu^m|^2
  -2\varepsilon\operatorname{Re}(\mu^m)
  +\varepsilon^2\\
&=1-2\varepsilon\operatorname{Re}(\mu^m)
  +\varepsilon^2\\
&>1.
\end{aligned}
\]

Thus
\[
\rho(T)\geq|\mu^m-\varepsilon|>1.
\]
\end{frame}

\begin{frame}{Proposition 5.1: Strict Dominance (IV)}
\small

We have shown that $\rho(T)>1$.

\begin{block}{Contradiction}But $0\leq T\leq M^m$, so monotonicity gives
\[
\rho(T)\leq\rho(M^m)=1.
\]
\end{block}


Thus $1$ is the only eigenvalue of $M$
with modulus one.

Since $\rho(M)=1$, such an eigenvalue
must exist. Hence $1$ is an eigenvalue
of $M$, so $r=\rho(A)>0$ is an eigenvalue
of $A$.

Every other eigenvalue $\lambda$ of $A$
therefore satisfies $|\lambda|<r$.
\end{frame}

\begin{frame}{Proposition 5.2: Positive Eigenvectors (I)}
\small

By Proposition 5.1, $r=\rho(A)>0$ is a real
eigenvalue. Choose a nonzero real eigenvector $x$:
\[
Ax=rx.
\]

Define $y=|x|$, taking absolute values
componentwise. Then $y\geq 0$ and $y\neq 0$.

Since $A$ has strictly positive entries,
the triangle inequality gives
\[
Ay\geq |Ax|=r|x|=ry.
\]

\begin{block}{Goal}
We want to show that equality holds:
\[
Ay=ry.
\]
This would give us a nonnegative eigenvector.
\end{block}
\end{frame}

\begin{frame}{Proposition 5.2: Positive Eigenvectors (II)}
\small

Suppose $Ay\geq ry$ is strict in at least
one coordinate.

Define
\[
d=Ay-ry\geq0,\qquad d\neq0.
\]

Because $A>0$ and $d$ is nonnegative
and nonzero, every component of $Ad$
is strictly positive:
\[
Ad>0.
\]

Multiplying $Ay=ry+d$ by $A$ gives
\[
A^2y=rAy+Ad>rAy.
\]

Now set $v=Ay>0$. Consequently,
\[
\boxed{Av>rv.}
\]
\end{frame}

\begin{frame}{Proposition 5.2: Positive Eigenvectors (III)}
\small

We have $v>0$ and $Av>rv$.

Define
\[
q=\min_i\frac{(Av)_i}{v_i}.
\]

Since the inequality is strict in every
coordinate,
\[
q>r,\qquad Av\geq qv.
\]


Repeated multiplication by $A$ gives
\[
A^kv\geq q^kv.
\]
\end{frame}

\begin{frame}{Proposition 5.2: Positive Eigenvectors (IV)}
\small

Recall that $v>0$, $q>r$, and $A^kv\geq q^kv$.

Taking infinity norms,
\[
\|A^k\|_\infty
\geq
\frac{\|A^kv\|_\infty}{\|v\|_\infty}
\geq q^k.
\]

By Gelfand's formula,
\[
\rho(A)=\lim_{k\to\infty}
\|A^k\|_\infty^{1/k}\geq q>r.
\]

This contradicts $\rho(A)=r$.
\end{frame}

\begin{frame}{Proposition 5.2: Positive Eigenvectors (V)}
\small

The contradiction shows that our
original inequality cannot be strict.
Therefore,
\[
Ay=ry.
\]

Since $A>0$ and $y\geq0$ is nonzero,
we know $Ay>0$.

Because $r>0$, the equation $Ay=ry$
implies $y>0$.

\begin{block}{Positive Right and Left Eigenvectors}
We have obtained a positive right
eigenvector $y>0$ satisfying $Ay=ry$.

Since $A^T>0$ and $\rho(A^T)=r$,
the same argument gives $w>0$ with
\[
A^Tw=rw
\quad\Longleftrightarrow\quad
w^TA=rw^T.
\]
\end{block}
\end{frame}

\begin{frame}{Proposition 5.3: Geometric Simplicity (I)}
\small

By Proposition 5.2, there exists $x>0$
such that $Ax=rx$.

Let $y$ be another real eigenvector
corresponding to $r$, so $Ay=ry$.

Define
\[
t=\max_i\frac{y_i}{x_i},
\qquad z=tx-y.
\]

By construction, $z\geq0$, and at least
one component of $z$ is zero.

Since $Ax=rx$ and $Ay=ry$,
\[
Az=tAx-Ay=r(tx-y)=rz.
\]
\end{frame}

\begin{frame}{Proposition 5.3: Geometric Simplicity (II)}
\small

Recall that $z\geq0$, at least one component of $z$ is zero, and $Az=rz$.

If $z\neq0$, then $A>0$ implies $Az>0$.
But $rz$ has a zero component,
giving a contradiction.

\begin{block}{Conclusion}
Thus $z=0$, so $y=tx$.
The result extends to complex eigenvectors
by considering real and imaginary parts.

Therefore,
\[
\boxed{\dim\ker(A-rI)=1.}
\]
\end{block}
\end{frame}

\begin{frame}{Proposition 5.3: Algebraic Simplicity (I)}
\small

By Proposition 5.2, there exist positive
right and left eigenvectors $x,w>0$ such that
\[
Ax=rx,\qquad A^Tw=rw.
\]

Suppose, for contradiction, that $r$
has algebraic multiplicity greater than one.

We already proved that its eigenspace is
one-dimensional:
\[
\dim\ker(A-rI)=1.
\]

Therefore, there exists a generalized
eigenvector $u$ satisfying
\[
(A-rI)u=x.
\]

\begin{block}{Key Idea}
We will multiply this equation by the
positive left eigenvector $w^T$ to
obtain a contradiction.
\end{block}
\end{frame}

\begin{frame}{Proposition 5.3: Algebraic Simplicity (II)}
\small

From the previous slide, we have
\[
(A-rI)u=x.
\]

Multiplying by the left eigenvector $w^T$ gives
\[
w^T(A-rI)u=w^Tx.
\]

Since $w^TA=rw^T$, we know
\[
w^T(A-rI)=w^TA-rw^T=0.
\]

\begin{block}{Contradiction}
\[
\boxed{0=w^Tx>0.}
\]

This is impossible because $w,x>0$,
so their dot product is strictly positive.

Thus, $r$ must be algebraically simple.
\end{block}
\end{frame}
\end{document}
